\chapter*{Glossary}
\addcontentsline{toc}{chapter}{Glossary}
\markboth{Glossary}{}

\begin{description}[leftmargin=0pt,style=unboxed,itemsep=6pt]

\item[Adversarial perturbation] A small, often visually diffuse change to
an image that flips a classifier's prediction without altering the
semantic content a human would recognise. Distinguishing genuine
counterfactuals from adversarial perturbations is the central validation
problem of this project.

\item[Background lock] A hard spatial constraint, adapted from RePaint and
Blended Latent Diffusion, in which the latent outside a permitted edit
mask is overwritten at every denoising step with the original image's
forward-noised latent, guaranteeing that region cannot be repainted.

\item[BDC (Breast Density Contrast)] A mammography diagnostic metric ---
mean intensity of the segmented dense breast tissue --- purpose-built to
detect a diffuse global brightness/density change large enough to move a
classifier without altering a lesion, the mammography signature of a
failure mode first identified on RSNA.

\item[CBIS-DDSM] The Curated Breast Imaging Subset of the Digital Database
for Screening Mammography: a public mammography dataset of mass and
calcification lesion cases labelled malignant or benign, used as the
second imaging domain in Chapter~\ref{ch:breast}. Unlike RSNA, it contains
no normal (lesion-free) images.

\item[Counterfactual explanation] For a query image $x$ classified by $f$,
an image $x_c$ that $f$ assigns to a different target class while
remaining realistic and preserving the identity of $x$. The difference
between $x$ and $x_c$ constitutes the explanation.

\item[Counterfactual Validity (CV)] The fraction of generated
counterfactuals for which the classifier's score both moves by more than
a margin $\tau$ and crosses the decision boundary in the intended
direction.

\item[CTR (Cardio-Thoracic Ratio)] The ratio of the heart's horizontal
span to the thorax's, computed here from segmentation masks. A ratio
above $0.5$ is the standard clinical cardiomegaly threshold.

\item[DDIM] Denoising Diffusion Implicit Model: a deterministic,
non-Markovian sampling procedure for diffusion models that permits
partial forward-then-reverse trajectories, and therefore image editing
rather than only synthesis from noise.

\item[DomainProfile / domain registry] A dataclass and module registry
(\file{domains/}) binding a dataset, its classifier target and its
clinical metrics together, so that adding a new imaging domain --- as in
Chapter~\ref{ch:breast} --- is one new module plus one configuration line
rather than a change threaded through the guidance and evaluation code.

\item[FID (Fr\'echet Inception Distance)] A distributional realism
measure comparing the means and covariances of deep features extracted
from real and generated image sets. Lower is better. Computed here in
TorchXRayVision DenseNet feature space rather than ImageNet Inception
space.

\item[FOP (Foreign-Object Preservation)] A metric measuring whether
medical hardware --- tubes, leads, pacemakers --- survives counterfactual
editing. Requires a dedicated device detector, and is not computable
offline in this project.

\item[GAE (Generalized Advantage Estimation)] A variance-reduction
technique for policy-gradient methods that interpolates between one-step
and Monte-Carlo advantage estimates via a parameter $\lambda$.

\item[Grad-CAM] Gradient-weighted Class Activation Mapping: a technique
producing class-discriminative localisation maps by weighting
convolutional feature maps with the gradients flowing into them. Used
here both as a diagnostic and as a spatial constraint.

\item[Guidance] The injection of a classifier's gradient into a diffusion
model's sampling trajectory, steering generation toward a target class
without retraining the generator.

\item[Held-out validation] Re-scoring generated counterfactuals with
classifiers that never participated in the guidance loop, to distinguish
genuine semantic change from adversarial perturbation against the
guidance model.

\item[Identity Preservation (IP)] A model-free substitute for FOP used
here, combining MS-SSIM (global anatomical similarity) with the
correlation of high-pass-filtered images (survival of thin, high-contrast
structures).

\item[LHR (Lesion Hit Rate)] The share of an edit's absolute-difference
mass falling inside the radiologist's annotated lesion box, divided by the
box's own area as a null baseline (so ``lift'' rather than raw hit rate
carries information). The only metric in this project anchored to
ground-truth disease location rather than to a classifier or a structural
proxy; introduced in Chapter~\ref{ch:integrity}.

\item[Latent diffusion] A diffusion model operating in the compressed
latent space of a pretrained autoencoder rather than in pixel space,
reducing computational cost by orders of magnitude.

\item[LLR (Lesion Load Ratio)] The mammography analogue of CTR: the ratio
of segmented lesion area to segmented breast area, expected to rise under
a benign$\to$malignant edit and fall under the reverse.

\item[LoRA (Low-Rank Adaptation)] A parameter-efficient fine-tuning
technique that freezes pretrained weights and trains small injected
low-rank matrix pairs, drastically reducing the memory and compute
required to adapt a large model.

\item[MOP (Marker/clip Preservation)] The mammography counterpart of FOP:
the fraction of images containing a surgical clip or skin marker in which
that object survives editing. Computable offline via a model-free
saturated-blob detector, unlike FOP.

\item[MS-SSIM] Multi-Scale Structural Similarity Index Measure: an image
similarity metric evaluating luminance, contrast and structure across
multiple spatial scales.

\item[MSI (Margin Spiculation Index)] A mammography clinical proxy metric:
the normalised boundary complexity of the segmented lesion's largest
component, higher for spiculated (malignancy-associated) margins than for
smoothly circumscribed ones. Replaces SCP in the mammography domain.

\item[PPO (Proximal Policy Optimization)] A policy-gradient reinforcement
learning algorithm that constrains each update with a clipped surrogate
objective, trading some theoretical rigour for stability and sample
efficiency.

\item[Reward shaping] The addition of auxiliary terms to a reinforcement
learning reward to guide learning toward desired behaviour. A recurring
source of bugs in this project, since the terms' relative magnitudes ---
not their presence --- determine what the agent actually optimises.

\item[SCP (Costophrenic-recess score)] A proxy for pleural-effusion-related
blunting of the costophrenic angle, approximated here from an object
detector's feature norm in a lower-lobe crop.

\item[SLURM] Simple Linux Utility for Resource Management: the batch
scheduler allocating compute nodes on the HPC cluster used for all
experiments.

\item[Transfer ratio] The fraction of the guidance classifier's score
shift that an independent held-out classifier reproduces on the same
counterfactual. Low values indicate an adversarial rather than semantic
edit.

\item[VAE (Variational Autoencoder)] The encoder-decoder pair mapping
images to and from the compressed latent space in which the diffusion
model operates. Frozen throughout this project.

\end{description}

\clearpage
